% LIGHT_Kr system map for v0_1: energy flow through pure krypton at LIGHT's
% reference point, on the kinetic EEDF, every band at its default. Every
% number comes from system_map_v0_1_numbers.tex, which
% make_system_map_v0_1.py writes from the model. Build from LIGHT_Kr/docs:
%   python3 make_system_map_v0_1.py && latexmk -pdf system_map_v0_1.tex
% Palette and styles are LIGHT's maps'.
\documentclass[11pt]{article}
\usepackage[a4paper,landscape,margin=0.9cm,top=0.7cm,bottom=0.6cm]{geometry}
\usepackage[T1]{fontenc}
\usepackage[scaled=0.95]{helvet}
\renewcommand{\familydefault}{\sfdefault}
\usepackage{amsmath,amssymb}
\usepackage{sansmath}
\usepackage{xcolor}
\usepackage{tikz}
\usetikzlibrary{arrows.meta,calc,positioning}
\pagestyle{empty}
\setlength{\parindent}{0pt}

% LIGHT_Kr v0_1 map numbers. GENERATED by make_system_map_v0_1.py
% from the model (0.1.0-kr). Do not edit. Rerun the script instead.
\providecommand\wmdef[2]{\expandafter\def\csname wm@#1\endcsname{#2}}
\wmdef{e.P}{100.0}
\wmdef{e.pen}{2.6}
\wmdef{e.pen2}{0.0}
\wmdef{e.G1}{19.7}
\wmdef{e.C1}{1.1}
\wmdef{e.C2}{0.0}
\wmdef{e.76}{77.9}
\wmdef{e.stepS}{0.0}
\wmdef{e.stepP}{0.0}
\wmdef{e.G6}{2.7}
\wmdef{e.dU}{0.0}
\wmdef{e.cap}{1.1}
\wmdef{C1.kr}{1.1}
\wmdef{C1.heat}{0.0}
\wmdef{kpp.back}{6.0}
\wmdef{KPS.heat}{78.1}
\wmdef{PRAD.hv}{0.6}
\wmdef{PRAD.heat}{0.0}
\wmdef{kpp.stepP}{0.0}
\wmdef{kpp.left}{0.0}
\wmdef{krs.G2}{8.9}
\wmdef{krs.AK8}{16.8}
\wmdef{krs.stepS}{0.0}
\wmdef{krs.left}{0.0}
\wmdef{G2.ion}{6.3}
\wmdef{C2.ion}{0.0}
\wmdef{stepS.ion}{0.0}
\wmdef{stepP.ion}{0.0}
\wmdef{AK8.heat}{1.2}
\wmdef{k2s.RAD}{15.3}
\wmdef{k2s.pen}{0.0}
\wmdef{k2s.left}{0.2}
\wmdef{RAD.hv}{14.1}
\wmdef{RAD.heat}{1.3}
\wmdef{pen2.ion}{0.0}
\wmdef{K15.heat}{0.6}
\wmdef{R11p.in}{6.4}
\wmdef{R11p.kr}{5.7}
\wmdef{R11p.heat}{0.8}
\wmdef{ions.left}{-0.7}
\wmdef{out.147}{14.1}
\wmdef{out.IR}{0.6}
\wmdef{out.heat}{83.1}
\wmdef{out.elastic}{2.7}
\wmdef{out.ions}{-0.7}
\wmdef{out.other}{0.2}
\wmdef{version}{0.1.0-kr}
\wmdef{md5}{9d36580be5ba}
\wmdef{p}{1}
\wmdef{Tg0}{300}
\wmdef{P}{300}
\wmdef{pulse}{1}
\wmdef{tend}{2}
\wmdef{ne}{10^{15}}
\wmdef{N}{2.45\times10^{19}}
\wmdef{Edep}{0.30}
\wmdef{Te}{2.29}
\wmdef{TeMax}{2.56}
\wmdef{enLo}{1.3}
\wmdef{enHi}{15.8}
\wmdef{kinN}{381}
\wmdef{kinOn}{276}
\wmdef{kinRatioLo}{1.0001}
\wmdef{kinRatioHi}{1.0002}
\wmdef{kinGap}{7\times10^{-6}}
\wmdef{closure}{4\times10^{-15}}
\wmdef{TgEnd}{808}
\wmdef{TgMax}{808}
\wmdef{pEnd}{2.69}
\wmdef{nePeak}{1.00\times10^{15}}
\wmdef{neEnd}{4.43\times10^{12}}
\wmdef{krsFormed}{1.14\times10^{18}}
\wmdef{kppFormed}{1.11\times10^{18}}
\wmdef{k2sFormed}{3.16\times10^{16}}
\wmdef{brs.76}{95.8}
\wmdef{brs.AK8}{2.8}
\wmdef{brs.G2}{1.5}
\wmdef{brp.KPS}{95.4}
\wmdef{brp.76r}{3.8}
\wmdef{brp.PRAD}{0.8}
\wmdef{brk2.RAD}{98.7}
\wmdef{k2s.leftpc}{1.3}
\wmdef{cyc.perK2}{35}
\wmdef{cyc.e}{77.9}
\wmdef{cyc.net}{1.05\times10^{18}}
\wmdef{ion.G2}{100}
\wmdef{ion.C2}{0.013}
\wmdef{k.G2}{5.0\times10^{-10}}
\wmdef{k.AK8}{1.7\times10^{-32}}
\wmdef{tpow.AK8}{-0.64}
\wmdef{tauRAD}{270}
\wmdef{hv.RAD}{8.43}
\wmdef{E.kr2s}{9.20}
\wmdef{E.kr2p}{12.846}
\wmdef{k.K15}{2.4\times10^{-31}}
\wmdef{tpow.K15}{-0.75}
\wmdef{k.R11p}{1.6\times10^{-6}}
\wmdef{pow.R11p}{0.55}
\wmdef{k.KPS}{1.7\times10^{-10}}
\wmdef{tauPRAD}{29.0}
\wmdef{E.PP}{11.30}
\wmdef{E.S}{9.915}
\wmdef{gap5p}{1.385}
\wmdef{pen.e}{5.834}
\wmdef{sig.76}{1.1\times10^{-14}}
\wmdef{nodes}{18}
\wmdef{notes.krs}{20.8}
\wmdef{notes.krp}{6.3}
\wmdef{notes.k2s}{14.2}
\wmdef{om.kannari}{15.23}
\wmdef{om.nocycle}{51.06}
\wmdef{om.twot}{15.05}
\wmdef{om.twot.te}{2.26}
\wmdef{bd.k_5s5p.def}{1}
\wmdef{bd.k_5s5p.lo}{0.5}
\wmdef{bd.k_5s5p.hi}{1.4}
\wmdef{bd.k_5s5p.lo147}{18.58}
\wmdef{bd.k_5s5p.hi147}{12.23}
\wmdef{bd.kps_kr.def}{1.7\times10^{-10}}
\wmdef{bd.kps_kr.lo}{3.0\times10^{-11}}
\wmdef{bd.kps_kr.hi}{1.7\times10^{-10}}
\wmdef{bd.kps_kr.lo147}{15.59}
\wmdef{bd.kps_kr.hi147}{14.06}
\wmdef{bd.kr_pen.def}{5.0\times10^{-10}}
\wmdef{bd.kr_pen.lo}{1.5\times10^{-10}}
\wmdef{bd.kr_pen.hi}{5.0\times10^{-10}}
\wmdef{bd.kr_pen.lo147}{18.37}
\wmdef{bd.kr_pen.hi147}{14.06}
\wmdef{bd.kr2_form.def}{1.7\times10^{-32}}
\wmdef{bd.kr2_form.lo}{1.7\times10^{-32}}
\wmdef{bd.kr2_form.hi}{5.5\times10^{-32}}
\wmdef{bd.kr2_form.lo147}{14.06}
\wmdef{bd.kr2_form.hi147}{34.70}
\wmdef{bd.kr2_pen.def}{0}
\wmdef{bd.kr2_pen.lo}{0}
\wmdef{bd.kr2_pen.hi}{5.0\times10^{-10}}
\wmdef{bd.kr2_pen.lo147}{14.06}
\wmdef{bd.kr2_pen.hi147}{7.67}
\wmdef{bd.kstep_s.def}{0}
\wmdef{bd.kstep_s.lo}{0}
\wmdef{bd.kstep_s.hi}{1.5\times10^{-7}}
\wmdef{bd.kstep_s.lo147}{14.06}
\wmdef{bd.kstep_s.hi147}{8.48}
\wmdef{bd.kstep_p.def}{0}
\wmdef{bd.kstep_p.lo}{0}
\wmdef{bd.kstep_p.hi}{2.0\times10^{-7}}
\wmdef{bd.kstep_p.lo147}{14.06}
\wmdef{bd.kstep_p.hi147}{11.27}
\wmdef{bd.n}{7}
\wmdef{bd.nzero}{3}
\def\wmBands{\el{} + \Krs{} \ra \el{} + \Krpp{}, a scale on the measured cross section & $1$ & $0.5$ to $1.4$ & 18.58 to 12.23 & 20.06 to 12.98 & 0.58 to 0.67 & 782 to 819\\
\Krpp{} + Kr \ra \Krs{} + Kr & $1.7\times10^{-10}$ & $3.0\times10^{-11}$ to $1.7\times10^{-10}$\unit{cm^3\,s^{-1}} & 15.59 to 14.06 & 16.56 to 15.05 & 3.42 to 0.65 & 783 to 808\\
\Krs{} + \Krs{} \ra \Krp{} + Kr + \el{} & $5.0\times10^{-10}$ & $1.5\times10^{-10}$ to $5.0\times10^{-10}$\unit{cm^3\,s^{-1}} & 18.37 to 14.06 & 18.86 to 15.05 & 0.62 to 0.65 & 782 to 808\\
\Krs{} + 2Kr \ra \Krts{} + Kr & $1.7\times10^{-32}$ & $1.7\times10^{-32}$ to $5.5\times10^{-32}$\unit{cm^6\,s^{-1}} & 14.06 to 34.70 & 15.05 to 35.29 & 0.65 to 0.45 & 808 to 687\\
\Krts{} + \Krts{} \ra \Krtp{} + 2Kr + \el{} & $0$ & $0$ to $5.0\times10^{-10}$\unit{cm^3\,s^{-1}} & 14.06 to 7.67 & 15.05 to 8.18 & 0.65 to 0.69 & 808 to 846\\
\el{} + \Krs{} \ra \Krp{} + 2\el{}, pre-exponential & $0$ & $0$ to $1.5\times10^{-7}$\unit{cm^3\,s^{-1}} & 14.06 to 8.48 & 15.05 to 9.90 & 0.65 to 0.64 & 808 to 842\\
\el{} + \Krpp{} \ra \Krp{} + 2\el{}, pre-exponential & $0$ & $0$ to $2.0\times10^{-7}$\unit{cm^3\,s^{-1}} & 14.06 to 11.27 & 15.05 to 13.01 & 0.65 to 0.66 & 808 to 825\\}

\makeatletter
\newcommand\wm[1]{\ifcsname wm@#1\endcsname\csname wm@#1\endcsname
  \else\wm@missing{#1}\fi}
\newcommand\wm@missing[1]{\errmessage{system_map_v0_1_numbers.tex has no
  number '#1'. Rerun make_system_map_v0_1.py}}
\makeatother

% LIGHT's palette: electrons, Kr*, Kr** and its near-IR light, Kr2* and
% 147 nm, heat, ions (stored, a neutral grey)
\definecolor{cE}{HTML}{2A78D6}
\definecolor{cX}{HTML}{EDA100}
\definecolor{cP}{HTML}{C46A00}
\definecolor{cV}{HTML}{A63A83}
\definecolor{cH}{HTML}{E34948}
\definecolor{cI}{HTML}{52514E}
\definecolor{ink}{HTML}{0B0B0B}
\definecolor{ink2}{HTML}{52514E}
\definecolor{muted}{HTML}{898781}

\newcommand{\unit}[1]{\ensuremath{\,\mathrm{#1}}}
\newcommand{\el}{\ensuremath{\mathrm{e}^-}}
\newcommand{\Krs}{Kr\ensuremath{^*}}
\newcommand{\Krpp}{Kr\ensuremath{^{**}}}
\newcommand{\Krp}{Kr\ensuremath{^+}}
\newcommand{\Krtp}{Kr\ensuremath{_2^+}}
\newcommand{\Krts}{Kr\ensuremath{_2^*}}
\newcommand{\ra}{\ensuremath{\to}\ }
\newcommand{\hnu}{\ensuremath{h\nu}}
\newcommand{\Edep}{\ensuremath{E_\mathrm{dep}}}
\newcommand{\pc}[1]{\wm{#1}\,\%}
\newcommand{\cm}[1]{\ensuremath{\,\mathrm{cm^{#1}s^{-1}}}}

\tikzset{
  band/.style={line width={max(0.6pt,(#1)*0.03mm)},
    -{Triangle[width=2.2pt 1.8,length=2pt 1.0]}},
  flat/.style={line width={max(0.6pt,(#1)*0.03mm)}},
  casing/.style={line width={max(0.6pt,(#1)*0.03mm)+2.4pt}, white},
  thin arrow/.style={line width=0.6pt,
    -{Triangle[width=2.2pt 1.8,length=2pt 1.0]}},
  blab/.style={font=\scriptsize, text=ink, inner sep=0.8pt, fill=white,
    fill opacity=0.85, text opacity=1},
  nohyph/.style={execute at begin node={\hyphenpenalty=10000
    \exhyphenpenalty=10000}},
  proc/.style={nohyph, draw=ink2!55, fill=white, rounded corners=2pt,
    text width=2.6cm, align=left, font=\scriptsize, inner sep=3pt,
    text=ink},
  spec/.style={nohyph, draw=#1, fill=#1!12, rounded corners=5pt,
    align=center, font=\scriptsize, text=ink, inner sep=3pt,
    text width=2.3cm, line width=0.8pt},
  inbox/.style={nohyph, draw=ink2!55, fill=white, rounded corners=2pt,
    text width=2.5cm, align=left, font=\scriptsize, inner sep=4pt,
    text=ink},
  tag/.style={draw=#1, fill=#1!14, rounded corners=1.5pt,
    font=\tiny, text=ink, inner sep=1.6pt},
  hdr/.style={font=\footnotesize\bfseries, text=ink2},
  note/.style={nohyph, font=\tiny, text=ink2},
}

\begin{document}
\sansmath

{\large\bfseries LIGHT\_Kr system map, v0\_1: energy flow through pure
krypton}\par
\vspace{2pt}
{\footnotesize Reference point: \wm{p}\unit{atm} of Kr at
\wm{Tg0}\unit{K}, $P=\wm{P}\unit{kW\,cm^{-3}}$ into the electrons for
\wm{pulse}\unit{\mu s} so that $\Edep=\wm{Edep}\unit{J\,cm^{-3}}$, $n_e$
seeded at $\wm{ne}\unit{cm^{-3}}$ and then solved, every band at its
default, run to \wm{tend}\unit{\mu s} on the kinetic EEDF. Band width is
proportional to the energy flow. Numbers are percent of \Edep, rounded to
0.1.}\par
\vspace{2pt}

\begin{tikzpicture}[x=1cm,y=1cm]
% ================================================================ frames
\node[hdr] at (1.35,15.7) {INPUT};
\node[hdr] at (25.8,15.7) {OUTPUT at $t=\wm{tend}\unit{\mu s}$};
\node[note] at (25.8,15.32) {percent of \Edep, sums to 100};
\draw[draw=ink2!35, fill=black!2.5, rounded corners=6pt]
  (2.95,0.25) rectangle (23.55,15.45);
\node[hdr, anchor=west] at (3.15,15.08)
  {UNIFORM INITIAL GAS\quad 0-D, constant volume};

% ================================================================ inputs
\node[inbox, line width=0.9pt, draw=cE] (pow) at (1.35,12.6) {%
  \textbf{Power in}\\
  $P=\wm{P}\unit{kW\,cm^{-3}}$\\ for \wm{pulse}\unit{\mu s}, uniform,\\
  into the electrons\\
  $\Edep=\wm{Edep}\unit{J\,cm^{-3}}$\\
  \textcolor{ink2}{the only input}};
\node[inbox] (seed) at (1.35,8.8) {%
  \textbf{$n_e$ seed}\\
  $\wm{ne}\unit{cm^{-3}}$ of \el{} and \Krp{} at $t=0$, then solved\\
  $\wm{neEnd}$ at the end};
\node[inbox] (fill) at (1.35,4.6) {%
  \textbf{Initial gas fill}\\
  pure Kr\\ \wm{p}\unit{atm}, \wm{Tg0}\unit{K}\\
  $N=\wm{N}$\\ \hspace*{1.2em}$\unit{cm^{-3}}$};

% ================================================================ electron gas
\node[draw=cE, fill=cE!10, rounded corners=5pt, line width=0.9pt,
  minimum width=1.6cm, minimum height=11.2cm, anchor=north]
  (eg) at (3.95,14.6) {};
\node[align=center, font=\scriptsize, text=ink, text width=1.5cm,
  anchor=north] at ([yshift=-4pt]eg.north) {\textbf{electron}\\
  \textbf{gas}\\[4pt] kinetic\\ EEDF\\ \textcolor{ink2}{two-term}\\
  \textcolor{ink2}{Boltzmann,}\\ \textcolor{ink2}{with e-e}\\[5pt]
  $T_e$ \wm{Te}\unit{eV}\\ \textcolor{ink2}{pulse mean}\\[5pt]
  $E/N$ \wm{enLo}\\ to \wm{enHi}\unit{Td}\\
  \textcolor{ink2}{in the pulse}\\[5pt] $U_e=\tfrac32 n_e T_e$};
\node[align=center, note, text width=1.4cm, anchor=south]
  at ([yshift=3pt]eg.south) {in: power \wm{e.P}\\ and Penning
  \wm{e.pen}\\[3pt] captured \el\\ \wm{e.cap} to heat};

% ================================================================ e-impact rows
\node[note, anchor=south west, text width=5.4cm, align=left]
  at (5.3,14.1) {$Q=k\,n_e[\mathrm{X}]$, $k=\langle\sigma v\rangle$ over
  the kinetic EEDF, on Biagi v7.1's cross sections and Jung et al.'s};
\node[proc, text width=2.35cm] (G1) at (6.6,13.35) {%
  \textbf{G1}\quad \el{} + Kr\\ \ra \el{} + \Krs(5s)\\
  \textcolor{ink2}{threshold \wm{E.S}\unit{eV}}};
\node[proc, text width=2.35cm] (C1) at (6.6,11.2) {%
  \textbf{C1}\quad \el{} + Kr\\ \ra Kr(5p, 4d, \dots)\\
  \textcolor{ink2}{cascades to \Krpp(5p)}};
\node[proc, text width=2.35cm] (K76) at (6.6,8.6) {%
  \textbf{K76, K76r}\\ \el{} + \Krs{} $\rightleftarrows$\\ \el{} +
  \Krpp\\ \textcolor{ink2}{measured $\sigma$:
  $\wm{sig.76}\unit{cm^2}$ at 6\unit{eV}}\\
  \textcolor{ink2}{\wm{gap5p}\unit{eV} a net cycle}};
\node[proc, text width=2.35cm] (C2) at (6.6,6.0) {%
  \textbf{C2}\quad \el{} + Kr\\ \ra \Krp{} + 2\el\\
  \textcolor{ink2}{threshold 13.996\unit{eV}}};
\node[proc, text width=2.35cm] (G6) at (6.6,4.35) {%
  \textbf{G6}\quad elastic\\ \el{} + Kr\\
  \textcolor{ink2}{$(2m/M)$ per collision}};

% ================================================================ Kr*, Kr**, R11p
\node[spec=cX] (KS) at (10.05,12.95) {%
  \textbf{\Krs(5s)} \wm{E.S}\unit{eV}\\
  formed $\wm{krsFormed}\unit{cm^{-3}}$\\
  {\tiny K76 \pc{brs.76}, AK8\_Kr \pc{brs.AK8}, G2 \pc{brs.G2}}};
\node[spec=cP] (KPP) at (10.05,8.6) {%
  \textbf{\Krpp(5p)} \wm{E.PP}\unit{eV}\\
  formed $\wm{kppFormed}\unit{cm^{-3}}$\\
  {\tiny KPS \pc{brp.KPS}, K76r \pc{brp.76r}, PRAD \pc{brp.PRAD}}};
\node[proc, text width=2.3cm] (R11) at (10.05,4.5) {%
  \textbf{R11p}\quad \Krtp{} + \el\\ \ra \Krpp{} + Kr\\
  \textcolor{ink2}{$\wm{k.R11p}\cm{3}$}\\
  \textcolor{ink2}{$\times\,(300\,\mathrm{K}/T_e)^{\wm{pow.R11p}}$}};

% ================================================================ Kr2*, G2, ions
\node[spec=cV, text width=2.5cm] (K2S) at (13.75,12.95) {%
  \textbf{\Krts}, \pc{brs.AK8} of \Krs\\
  $\wm{k2sFormed}\unit{cm^{-3}}$\\
  {\tiny AK8\_Kr: \Krs{} + 2Kr}\\
  {\tiny $\wm{k.AK8}\cm{6}\times(T_g/300)^{\wm{tpow.AK8}}$}};
\node[proc, text width=2.5cm] (G2) at (13.75,10.35) {%
  \textbf{G2}\quad \Krs{} + \Krs\\ \ra \Krp{} + Kr + \el\\
  \textcolor{ink2}{$\wm{k.G2}\cm{3}$, \wm{pen.e}\unit{eV} to \el}\\
  takes \pc{brs.G2} of \Krs};
\node[spec=cI, text width=2.5cm] (KI) at (13.75,6.0) {%
  \textbf{\Krp{}, \Krtp{} ions}\\ {\tiny K15: \Krp{} + 2Kr \ra \Krtp{} +
  Kr}\\ {\tiny $\wm{k.K15}\cm{6}\times(T_g/300)^{\wm{tpow.K15}}$}};

% ================================================================ Kr2* light
\node[proc, text width=2.5cm] (KR) at (17.55,12.95) {%
  \textbf{KR2\_RAD}\quad 147\unit{nm}\\ \Krts{} \ra 2Kr + \hnu\\
  \textcolor{ink2}{$\tau=\wm{tauRAD}\unit{ns}$}, \pc{brk2.RAD} of \Krts};
\node[note, draw=ink2!40, dashed, rounded corners=2pt, fill=white,
  text width=3.3cm, align=left, inner sep=3pt] (Z) at (17.9,10.35) {%
  \textbf{No measurement, 0 by default}, printed and not drawn:
  \textbf{PEN\_KR2} \Krts{} + \Krts{} \ra \Krtp{} + 2Kr + \el;
  \textbf{KSTEP\_S, KSTEP\_P} \el{} + \Krs, \Krpp{} \ra \Krp{} + 2\el.
  Each is a band (bottom left).};

% ================================================================ bands
% inputs into the gas
\draw[band=\wm{e.P}, cE] (pow.east) -- (pow.east -| eg.west)
  node[blab, midway, above=2.2pt] {\wm{e.P}};
\draw[thin arrow, ink2] (seed.east) -- (seed.east -| eg.west)
  node[blab, midway, above=1pt, font=\tiny] {seeds $n_e$};
\draw[thin arrow, ink2] (fill.east) -- (fill.east -| 2.95,0)
  node[blab, midway, above=1pt, font=\tiny] {fills};
% electron gas out to the electron-impact rows
\foreach \n/\k in {G1/e.G1, C1/e.C1, K76/e.76, C2/e.C2, G6/e.G6}{
  \draw[band=\wm{\k}, cE] (\n.west -| eg.east) -- (\n.west)
    node[blab, midway, above=1.5pt] {\wm{\k}};
}
% into Kr* and into the 5p rung
\draw[band=\wm{e.G1}, cX] (G1.east) -- (KS.west |- G1.east)
  node[blab, midway, above=1.5pt] {\wm{e.G1}};
\draw[band=\wm{C1.kr}, cP] (C1.south east) -- (KPP.north west)
  node[blab, midway, above right=0pt and 0pt] {\wm{C1.kr}};
\draw[band=\wm{e.76}, cP] (K76.east) -- (KPP.west)
  node[blab, midway, above=4pt] {\wm{e.76}};
% the rung hands its Kr* energy back, net of K76, K76r, KPS and PRAD
\draw[band=\wm{kpp.back}, cX] ([xshift=0.55cm]KPP.north) --
  ([xshift=0.55cm]KS.south)
  node[blab, midway, right=2pt] {\wm{kpp.back}};
% out of Kr*: to Kr2*, to Penning
\draw[band=\wm{krs.AK8}, cX] (KS.east) -- (K2S.west)
  node[blab, midway, above=2pt] {\wm{krs.AK8}};
\draw[band=\wm{krs.G2}, cX] (KS.south east) -- (G2.north west)
  node[blab, midway, above right=0pt and 0pt] {\wm{krs.G2}};
% Kr2* to its light, and the light to the ledger
\draw[band=\wm{k2s.RAD}, cV] (K2S.east) -- (KR.west)
  node[blab, midway, above=2pt] {\wm{k2s.RAD}};
\draw[band=\wm{RAD.hv}, cV] (KR.east) -- (KR.east -| 23.8,0)
  node[blab, pos=0.85, above=2pt] {\wm{RAD.hv}};
% ionization into the ions
\draw[band=\wm{C2.ion}, cI] (C2.east) -- (C2.east -| KI.west)
  node[blab, pos=0.9, above=1.5pt] {\wm{C2.ion}};
\draw[band=\wm{G2.ion}, cI] ([xshift=0.4cm]G2.south) --
  ([xshift=0.4cm]G2.south |- KI.north)
  node[blab, pos=0.85, right=3pt] {\wm{G2.ion}};
% recombination: the ions to R11p, R11p up into the 5p rung over C2's band
\draw[band=\wm{R11p.in}, cI] (KI.south west) -- (R11.east)
  node[blab, midway, below right=0pt and 0pt] {\wm{R11p.in}};
\draw[casing=\wm{R11p.kr}] (R11.north) -- (KPP.south);
\draw[band=\wm{R11p.kr}, cP] (R11.north) -- (KPP.south)
  node[blab, pos=0.5, left=2pt] {\wm{R11p.kr}};
% the rung's near-IR light and its relaxation heat, right to the ledger,
% over the ionization band of G2
\draw[casing=\wm{PRAD.hv}] ([yshift=5pt]KPP.east) -- ([yshift=5pt]KPP.east -| 23.8,0);
\draw[band=\wm{PRAD.hv}, cP] ([yshift=5pt]KPP.east) -- ([yshift=5pt]KPP.east -| 23.8,0)
  node[blab, pos=0.93, above=1.5pt] {\wm{PRAD.hv}};
\node[note, anchor=south west] at ([xshift=4pt, yshift=6pt]KPP.east)
  {PRAD, $\tau=\wm{tauPRAD}\unit{ns}$, 5p \ra 5s};
\coordinate (kpsA) at ([yshift=-6pt]KPP.east);
\draw[casing=\wm{KPS.heat}, rounded corners=6pt]
  (kpsA) -- (kpsA -| 21.9,0) -- (21.9,5.55) -- (23.8,5.55);
\draw[band=\wm{KPS.heat}, cH, rounded corners=6pt]
  (kpsA) -- (kpsA -| 21.9,0) -- (21.9,5.55) -- (23.8,5.55);
\node[blab, anchor=north] at (18.1,7.88) {\textbf{KPS}, $\wm{k.KPS}\cm{3}$:
  the \wm{gap5p}-eV gap to heat, \wm{KPS.heat}};
% Penning back to the electrons, down beside G2's ionization band and
% left under the 5p rung, over the rung's and R11p's bands
\draw[casing=\wm{e.pen}, rounded corners=5pt]
  ([xshift=-0.4cm]G2.south) -- ([xshift=-0.4cm]G2.south |- 0,7.15)
  -- (eg.east |- 0,7.15);
\draw[band=\wm{e.pen}, cE, rounded corners=5pt]
  ([xshift=-0.4cm]G2.south) -- ([xshift=-0.4cm]G2.south |- 0,7.15)
  -- (eg.east |- 0,7.15)
  node[blab, pos=0.72, below=2pt] {Penning returns \wm{e.pen}};

% ================================================================ loss tags
\node[tag=cH, anchor=north east] at ([yshift=-2pt]C1.south east)
  {heat \wm{C1.heat}};
\node[tag=cH, anchor=north east] at ([yshift=-2pt]G6.south east)
  {heat \wm{e.G6}, elastic};
\node[tag=cH, anchor=north east] at ([yshift=-2pt]R11.south east)
  {heat \wm{R11p.heat}};
\node[tag=cH, anchor=north east] at ([yshift=-2pt]KI.south east)
  {heat \wm{K15.heat}, K15};
\node[tag=cH, anchor=north east] at ([yshift=-2pt]K2S.south east)
  {heat \wm{AK8.heat}, AK8\_Kr};
\node[tag=cH, anchor=north east] at ([yshift=-2pt]KR.south east)
  {heat \wm{RAD.heat}};

% ================================================================ bands box
\node[draw=ink2!60, dashed, rounded corners=3pt, fill=white,
  align=left, font=\tiny, text=ink, anchor=south west,
  text width=11.4cm, inner sep=4pt] (BX) at (5.3,0.35) {%
  {\scriptsize\textbf{Bands.}} The 147-nm share of \Edep{} at each end of
  one band, the rest at their defaults, on the kinetic EEDF (here
  \wm{out.147}). The default end is in bold; \wm{bd.nzero} of the
  \wm{bd.n} bands have no measurement and a default of 0.\\[2pt]
  \begin{tabular}{@{}l@{\hspace{3pt}}r@{\hspace{9pt}}l@{\hspace{3pt}}r@{\hspace{9pt}}l@{\hspace{3pt}}r@{}}
  \el{} + \Krs{} \ra \Krpp{}, scale & \wm{bd.k_5s5p.lo147} to
    \wm{bd.k_5s5p.hi147}
    & \Krpp{} + Kr \ra \Krs{} & \wm{bd.kps_kr.lo147} to
    \textbf{\wm{bd.kps_kr.hi147}}
    & \Krs{} + \Krs{} & \wm{bd.kr_pen.lo147} to
    \textbf{\wm{bd.kr_pen.hi147}}\\
  \Krs{} + 2Kr & \textbf{\wm{bd.kr2_form.lo147}} to
    \wm{bd.kr2_form.hi147}
    & \Krts{} + \Krts{} & \textbf{\wm{bd.kr2_pen.lo147}} to
    \wm{bd.kr2_pen.hi147}
    & \el{} + \Krs{} \ra \Krp{} & \textbf{\wm{bd.kstep_s.lo147}} to
    \wm{bd.kstep_s.hi147}\\
  \el{} + \Krpp{} \ra \Krp{} & \textbf{\wm{bd.kstep_p.lo147}} to
    \wm{bd.kstep_p.hi147} & & & & \\
  \end{tabular}\\[2pt]
  The scale on the measured 5s \ra 5p cross section runs from
  \wm{bd.k_5s5p.lo} (the metastables alone) to \wm{bd.k_5s5p.hi} (the
  measurement's upper error), default \wm{bd.k_5s5p.def}.
  With the cycle off (scale 0) the share is \pc{om.nocycle}; with Kannari's
  fit in K76's place \pc{om.kannari}; on the two-temperature closure
  \pc{om.twot}.};

% ================================================================ legend
\begin{scope}[shift={(17.2,0.4)}, font=\tiny, text=ink]
  \node[anchor=west, font=\tiny\bfseries, text=ink2] at (0,2.95)
    {Band color is where the energy sits};
  \foreach \c/\t [count=\i] in {cE/electron gas, cX/\Krs{} (5s),
      cP/{\Krpp{} (5p) and near-IR light},
      cV/\Krts{} and 147-nm light, cI/\Krp{} and \Krtp{} ions,
      cH/heat}{
    \draw[\c, line width=1.2mm] (0,2.95-0.36*\i) -- ++(0.5,0);
    \node[anchor=west] at (0.56,2.95-0.36*\i) {\t};
  }
  \node[tag=cH, anchor=west] at (3.55,2.59) {heat};
  \node[anchor=west] at (4.15,2.59) {turns into heat};
  \draw[ink2, line width={10*0.03mm}] (3.55,2.23) -- ++(0.5,0);
  \node[anchor=west] at (4.1,2.23) {10\,\% band};
  \draw[ink2, line width={50*0.03mm}] (3.55,1.87) -- ++(0.5,0);
  \node[anchor=west] at (4.1,1.87) {50\,\% band};
\end{scope}

% ================================================================ ledger
\begin{scope}[font=\scriptsize, text=ink]
  \newcommand{\ledger}[5]{% y, color, label, key, bar options
    \node[anchor=south west, inner sep=1pt] at (23.85,#1) {#3};
    \node[anchor=south east, inner sep=1pt, font=\scriptsize\bfseries]
      at (27.85,#1) {$\wm{#4}$};
    \draw[ink2!30] (23.9,#1-0.1) -- (27.8,#1-0.1);
    \fill[#2, #5] (23.9,#1-0.1) rectangle
      ++({3.9*max(0,\wm{#4})/100},-0.14);
  }
  \ledger{12.95}{cV}{147-nm emission}{out.147}{}
  \ledger{9.0}{cP}{near-IR emission}{out.IR}{}
  \ledger{5.55}{cH}{reaction heat}{out.heat}{}
  \ledger{4.6}{cH}{elastic losses}{out.elastic}{opacity=0.55}
  \ledger{3.65}{cI}{stored in ions}{out.ions}{}
  \ledger{2.7}{muted}{other stores}{out.other}{}
  \node[note, anchor=north west, text width=3.8cm, align=left]
    at (23.85,3.5) {the seed's ions recombine};
  \node[note, anchor=north west, text width=3.8cm, align=left]
    at (23.85,2.45) {The KPS band and the heat tags inside the gas add up
    to the reaction heat. Ledger closure: $\wm{closure}$ of \Edep. Every
    node above conserves its flows (\wm{nodes} checks).};
\end{scope}
\node[note, anchor=north west, text width=2.6cm, align=left]
  at (0.05,2.9) {Generated from LIGHT\_Kr v0\_1, version \wm{version}, the
  default network with every band at its default. No number on this page
  is typed by hand.};
\end{tikzpicture}

\vspace{1pt}
{\tiny The outputs as flows, which overlap: to \Krs{} and \Krpp{}
by electron impact \pc{notes.krs}, to \Krp{} \pc{notes.krp}, to \Krts{}
\pc{notes.k2s} (formed times the 147-nm quantum), heat \pc{out.heat},
elastic losses \pc{out.elastic}. \textbf{The 5s
$\leftrightarrow$ 5p cycle.} The electrons pump \Krs{} up to the 5p rung
(K76, on Jung et al.'s measured cross sections) \wm{cyc.perK2} times for
each \Krts{} formed, and krypton relaxes the rung back (KPS), so the cycle
is drawn as its two net bands: of the \Krs{} formed (gross), \pc{brs.76}
is pumped up; the rung relaxes \pc{brp.KPS} of it on Kr, radiates
\pc{brp.PRAD} and the electrons return \pc{brp.76r}; the
\wm{gap5p}\unit{eV} of every net cycle, $\wm{cyc.net}\unit{cm^{-3}}$ of
them, takes \pc{cyc.e} of the input from the electrons, and KPS turns it
into heat. Penning ionization (G2) makes \wm{ion.G2}\,\% of the ions, and
they recombine onto the rung (R11p). \textbf{The kinetic EEDF:}
\wm{kinN} solutions, \wm{kinOn} in the pulse; over its second half the
field carries \wm{kinRatioLo} to \wm{kinRatioHi} of the input and the
kinetic and the model's energy balances agree to $\wm{kinGap}$ of it. At the
end $T_g=\wm{TgEnd}\unit{K}$, $p=\wm{pEnd}\unit{atm}$, $n_e=\wm{neEnd}\unit{cm^{-3}}$.
\textbf{Sources} of the rates printed: cross sections Biagi v7.1 via LXCat,
Pancheshnyi et al., Chem.\ Phys.\ 398, 148 (2012), doi:10.1016/j.chemphys.2011.04.020;
K76 Jung et al., Phys.\ Rev.\ A 73, 022722 (2006), doi:10.1103/PhysRevA.73.022722;
AK8\_Kr Ueda et al., Rev.\ Laser Eng.\ 12, 357 (1984), doi:10.2184/lsj.12.357, on Wieme and
Lenaerts' law, J.\ Chem.\ Phys.\ 74, 483 (1981), doi:10.1063/1.440855;
KR2\_RAD Morikawa et al., J.\ Chem.\ Phys.\ 91, 1469 (1989), doi:10.1063/1.457108;
G2 Leichner and Ericson, Phys.\ Rev.\ A 9, 251 (1974), doi:10.1103/PhysRevA.9.251,
as Kannari et al.\ print it, J.\ Appl.\ Phys.\ 57, 4309 (1985), doi:10.1063/1.334590;
KPS Gazeli et al., Phys.\ Plasmas 28, 043301 (2021), doi:10.1063/5.0041471;
PRAD Dzier\.z\k{e}ga et al., Phys.\ Rev.\ A 62, 022505 (2000), doi:10.1103/PhysRevA.62.022505;
R11p Shiu and Biondi, Phys.\ Rev.\ A 16, 1817 (1977), doi:10.1103/PhysRevA.16.1817;
K15 the mean of Bhattacharya, Phys.\ Rev.\ A 1, 1196 (1970), doi:10.1103/PhysRevA.1.1196,
Gr\"ossl et al., J.\ Chem.\ Phys.\ 74, 1728 (1981), doi:10.1063/1.441261, and
Smith et al., J.\ Phys.\ B 5, 2134 (1972), doi:10.1088/0022-3700/5/11/020.\par}

\end{document}
